\section{Gauge fixing and zero modes}\label{app:gauge-fixing-zero-modes}

This appendix records the gauge-fixing details used in Section \ref{sec:mode-solutions}. We first explain how temporal gauge is reached, then fix the residual spatial gauge freedom, then reduce the equations to the stream-function equation, and finally isolate the static configurations that do not belong to the propagating oscillator sector.

\subsection{Reaching temporal gauge}\label{subsec:app-reaching-temporal-gauge}

For a general off-shell configuration $A_{\mu}$, temporal gauge can be reached locally by choosing the gauge parameter

\begin{align}
\Lambda(t,r,\phi)&=-\int^{t}\mathrm{d}t'\,A_{t}(t',r,\phi).
\end{align}

This gives $A'_{t}=A_{t}+\partial_{t}\Lambda=0$. The spatial components transform as

\begin{align}
A'_{r}&=A_{r}-\int^{t}\mathrm{d}t'\,\partial_{r}A_{t}=o(r^{-1}), \\
A'_{\phi}&=A_{\phi}-\int^{t}\mathrm{d}t'\,\partial_{\phi}A_{t}=O(r^{0}).
\end{align}

Thus the temporal-gauge construction preserves the spatial falloff class introduced in Section \ref{sec:maxwell-fields-global-ads3}. The explicit normalizable representatives selected after residual gauge fixing obey the stronger radial decay $A_{r}=O(r^{-3})$.

\subsection{Residual gauge fixing}\label{subsec:app-residual-gauge-fixing}

After imposing $A_{t}=0$, the residual gauge parameter becomes time independent. The remaining condition used in the main text is

\begin{align}
\mathcal{C}[A]&=\partial_{r}A^{r}+\frac{1-r^{2}}{r(1+r^{2})}A^{r}+\partial_{\phi}A^{\phi}=0.
\end{align}

It is convenient to rewrite this condition as

\begin{align}
\partial_{r}\left(\frac{r}{1+r^{2}}A^{r}\right)+\partial_{\phi}\left(\frac{r}{1+r^{2}}A^{\phi}\right)&=0.
\end{align}

In this form, the constraint is solved locally by introducing a stream function $\Psi$ such that

\begin{align}
A^{r}&=\frac{1+r^{2}}{r}\partial_{\phi}\Psi, &
A^{\phi}&=-\frac{1+r^{2}}{r}\partial_{r}\Psi.
\end{align}

For a generic temporal-gauge configuration, the residual gauge parameter can be chosen to solve

\begin{align}
\left[(1+r^{2})\partial_{r}^{2}+\frac{1+r^{2}}{r}\partial_{r}+\frac{1}{r^{2}}\partial_{\phi}^{2}\right]\Lambda(r,\phi)&=-\mathcal{C}[A].
\end{align}

This removes the proper residual gauge freedom, up to zero modes that must be treated separately.

\subsection{Reduction of the equations of motion}\label{subsec:app-eom-reduction}

Once the temporal-gauge potential is parameterized by $\Psi$,

\begin{align}
A^{t}&=0, &
A^{r}&=\frac{1+r^{2}}{r}\partial_{\phi}\Psi, &
A^{\phi}&=-\frac{1+r^{2}}{r}\partial_{r}\Psi,
\end{align}

the temporal-gauge Maxwell equations reduce to

\begin{align}
\frac{1+r^{2}}{r}\partial_{\phi}(\mathcal{D}_{0}\Psi)&=0, \\
-\left(\frac{1+r^{2}}{r}\partial_{r}+2\right)(\mathcal{D}_{0}\Psi)&=0,
\end{align}

where

\begin{align}
\mathcal{D}_{0}\Psi&=(1+r^{2})\partial_{r}^{2}\Psi+\frac{1+3r^{2}}{r}\partial_{r}\Psi-\frac{1}{1+r^{2}}\partial_{t}^{2}\Psi+\frac{1}{r^{2}}\partial_{\phi}^{2}\Psi.
\end{align}

The second equation implies that $\mathcal{D}_{0}\Psi$ can differ from zero only by a term of the form $C(t)/(1+r^{2})$. This contribution is removable by the redundancy $\Psi\mapsto\Psi+q(t)$, so the propagating sector may be represented by the reduced equation $\mathcal{D}_{0}\Psi=0$ used in Section \ref{sec:mode-solutions}.

\subsection{Static zero modes}\label{subsec:app-static-modes}

The remaining zero-energy configurations are analyzed separately. For static modes, write

\begin{align}
A^{\mu}(r,\phi)&=e^{im\phi}f^{\mu}(r).
\end{align}

For the axisymmetric sector $m=0$, the static equations allow

\begin{align}
f^{r}&=h_{0}(r), &
f^{\phi}&=\frac{C_{1}}{r^{2}}+C_{2}\frac{\log(1+r^{2})}{r^{2}}.
\end{align}

The arbitrary $f^{r}$ branch is generated by the residual gauge transformation

\begin{align}
\Lambda_{0}(r)&=\int^{r}\frac{h_{0}(\rho)}{1+\rho^{2}}\,\mathrm{d}\rho.
\end{align}

For nonzero angular momentum, the general static solution can be written as

\begin{align}
f^{\phi}&=g_{m}(r), \\
f^{r}&=-\frac{ir^{2}(1+r^{2})}{m}\left(g_{m}'(r)+\frac{2}{r}g_{m}(r)\right).
\end{align}

It is convenient to rewrite this solution in terms of

\begin{align}
\lambda_{m}(r)&=\frac{r^{2}}{im}g_{m}(r),
\end{align}

for which

\begin{align}
f^{r}&=(1+r^{2})\lambda_{m}'(r), &
f^{\phi}&=\frac{im}{r^{2}}\lambda_{m}(r).
\end{align}

Thus the $m\neq0$ static sector is also locally generated by a residual gauge transformation, now with parameter $\Lambda=e^{im\phi}\lambda_{m}(r)$. Therefore all static zero modes in this appendix are residual gauge configurations. They do not contribute additional propagating degrees of freedom and are excluded from the positive-frequency oscillator basis used in the main text.
